\begin{abstract}
Group Relative Policy Optimization (GRPO) standardises each reward against the
statistics of a group of $G$ completions for the same prompt. When all
completions in a group receive the same verifiable reward, every advantage is
exactly zero and the group contributes no reward-driven gradient, even though
it was sampled and scored at full cost. We study this \emph{signal
starvation} directly. We define the Zero-Variance Fraction (ZVF), recompute
it from stored reward tensors, and find that on our main configuration
roughly 0.72--0.77 of each batch is starved, mostly because prompts are
already solved rather than because they are too hard. A simple signal model,
$S = p(1-p)(1-h_G(p))$ with $h_G(p) = p^G + (1-p)^G$, makes three
falsifiable predictions; two reproduce at scale, including a U-shape of ZVF
across a 368-run audit, and the third holds only directionally at toy scale.
Because ZVF cannot tell mastery from incapacity and collapses under a
$10^{-4}$ reward jitter, we promote the pairwise contrast density (PCD) as the
control signal. Finally, we audit an adaptive group-size controller: it ties
the best fixed recipe (+0.575 held-out gain at 186 rollouts), but 92.3\% of its
escalations fire on all-correct groups, where no group size can restore
contrast. Curriculum and re-baselining interventions return nulls. We
conclude that ZVF and PCD are useful monitoring diagnostics, not objectives
to optimise.
\end{abstract}

\begin{IEEEkeywords}
reinforcement learning, GRPO, large language models, verifiable rewards,
advantage estimation, group size
\end{IEEEkeywords}

\section{Introduction}

Reinforcement learning with verifiable rewards (RLVR) trains a language model
against a programmatic checker, such as exact match for mathematics or unit
tests for code, instead of a learned reward model~\cite{ouyang2022training}.
GRPO~\cite{shao2024deepseekmath,guo2025deepseekr1} has become the default
optimiser for this setting because it removes the value network of
PPO~\cite{schulman2017proximal}. It samples $G$ completions per prompt and
standardises their rewards within the group:
\begin{equation}
A_i = \frac{r_i - \mu_g}{\sigma_g + \varepsilon}, \quad
\mu_g = \tfrac{1}{G}\textstyle\sum_j r_j, \quad
\sigma_g^2 = \tfrac{1}{G}\textstyle\sum_j (r_j-\mu_g)^2 .
\label{eq:adv}
\end{equation}
If every completion in the group receives the same reward, $\sigma_g = 0$ and
every $A_i$ is zero. The KL term still acts, but the group-relative part of
the update, which is the only part driven by the task, vanishes. Standard
training logs (reward, loss, gradient norm, entropy, length) do not show how
often this happens: a rising reward curve looks the same whether a fifth or
four fifths of the batch is starved.

This paper measures that quantity and asks what, if anything, it tells a
practitioner to do. Our contributions are:
\begin{itemize}
\item a precise, recomputable definition of the Zero-Variance Fraction and
its measurement on real GRPO reward tensors;
\item evidence that starvation in our regime is \emph{easy-driven}, and that
ZVF aliases mastery with incapacity, which motivates the pairwise contrast
density (PCD) as a more robust signal;
\item a signal model for group size with three falsifiable predictions, and
the measurements that test them;
\item a counterfactual audit of an adaptive group-size controller, together
with null results for curriculum and re-baselining interventions.
\end{itemize}

\section{Related Work}

DAPO~\cite{yu2025dapo} addresses degenerate groups by dynamic sampling,
discarding and resampling groups whose rewards are uniform.
GSPO~\cite{qwen2025gspo} moves the importance ratio to the sequence level,
and Dr.\ GRPO~\cite{liu2025understanding} removes the length and
standard-deviation normalisations that it shows to be biased. Open-source
training stacks such as TRL~\cite{vonwerra2022trl}, veRL~\cite{sheng2024verl}
and OpenRLHF~\cite{hu2024openrlhf} implement these variants with different
defaults. Henderson et al.~\cite{henderson2018deep} showed how sensitive deep
RL conclusions are to seeds and implementation details, which shapes our
statistical protocol. Unlike this work on new estimators, we do not propose
an optimiser; we measure how much of the group-relative signal is available
in the first place, and test whether acting on that measurement helps.

\section{Method}

\subsection{Zero-Variance Fraction}
For a batch $\mathcal{B}_t$ of prompt groups at step $t$, each with $K$
sampled completions and rewards normalised to $[0,1]$,
\begin{equation}
\mathrm{ZVF}_t = \frac{1}{|\mathcal{B}_t|}\sum_{g\in\mathcal{B}_t}
\mathbb{1}\!\left[\widehat{\mathrm{Var}}(r_{g,1},\dots,r_{g,K}) \le \epsilon\right],
\end{equation}
with $\widehat{\mathrm{Var}}$ the unbiased sample variance and
$\epsilon = 10^{-6}$. Gradient utilisation is its complement,
$\mathrm{GU}_t = 1 - \mathrm{ZVF}_t$. This is an accounting identity about the
centred reward contrast, not a claim that the total gradient is zero:
objectives with KL, clipping or auxiliary losses still produce an update.
Sensitivity to $\epsilon$ is at most 0.4\%. The definition does not apply to
dense rewards, $K = 1$, or non-outcome rewards.

\subsection{Pairwise contrast density}
ZVF is unsigned: an all-wrong group and an all-correct group score
identically. We therefore also track the pairwise contrast density,
\begin{equation}
\mathrm{PCD} = \tfrac{G-1}{G}\,\mathbb{E}\left[p(1-p)\right],
\end{equation}
the expected fraction of within-group pairs whose rewards differ.

\subsection{A signal model for group size}
For a prompt with per-rollout success probability $p$ and $G$
conditionally independent Bernoulli rewards, the probability that the group
is degenerate is $h_G(p) = p^G + (1-p)^G$. We model usable signal per
prompt as
\begin{equation}
S(p,G) = p(1-p)\bigl(1 - h_G(p)\bigr).
\label{eq:signal}
\end{equation}
The first factor is the contrast a non-degenerate group can carry; the second
is the probability that any contrast survives centring. Both vanish at
$p \in \{0,1\}$ and peak at $p = \tfrac12$. Three consequences follow:
\textbf{(T1)} the expected ZVF indicator of a group equals $h_G(p)$;
\textbf{(T2)} $h_G$ is strictly decreasing in $G$, at rate
$-\ln\max(p,1-p)$; and \textbf{(T3)} with unnormalised group-centred
advantages, the mean absolute advantage is $2\hat p(1-\hat p)$, so gradient
magnitude should track $p(1-p)$.

\subsection{Experimental setup}
Training runs use LoRA~\cite{hu2022lora} adapters on Qwen-family models
from 0.5B to 32B parameters, together with Llama-3 models, against GSM8K
exact-match~\cite{cobbe2021training} and synthetic arithmetic verifiers. The
training loop logs per-step ZVF, gradient utilisation, entropy, completion
length and KL, and persists the raw per-group reward tensors so that every
diagnostic can be recomputed rather than read from a log line. The same
task, reward and decoding specification is dispatched to a managed training
API and to an open trainer; comparisons hold the stack fixed and vary one
factor. Held-out evaluation uses disjoint splits, and every gain is reported
with its $n$.

\section{Results}

\subsection{Starvation is the typical, easy-driven regime}
On the main configuration ZVF is approximately 0.72--0.77: about three
quarters of each rollout batch contributes no group-relative gradient. The
statistic is mechanically checked by recomputation: three stored GSM8K
group tensors give 0.130, 0.190 and 0.155 (pooled 0.1583), exactly matching
the logged values. A matched-budget panel of 2,560 rollouts per arm on
Qwen3-8B/GSM8K shows the mechanism. Late $G=2$ arms reach a training reward
of about 1.0 with ZVF between 0.75 and 1.0, while $G=16$ arms hold ZVF at or
below 0.25 at the same budget. When the policy solves the task, nearly every
group becomes all-correct. Starvation here is driven by ease, not
difficulty.

\subsection{ZVF is a diagnostic, not a score}
Four observations argue against optimising ZVF directly. First, it aliases
mastery with incapacity: in a 368-run audit across seven models, ZVF is
0.95--0.97 at both reward extremes and 0.25--0.29 at mid-range rewards of
0.35--0.50. Second, its association with final held-out outcome is weak
($\rho \approx 0.27$), although it correlates better with catastrophic
collapse (Spearman $\rho = 0.56$). Third, ZVF rises only 1--3 steps before a
reward plateau, so it is a warning signal rather than a forecast. Fourth, it
is fragile: adding micro-jitter $\epsilon \sim U(0, 10^{-4})$ to the rewards
collapses batch ZVF from 0.158 to 0.000, while PCD moves by only
$3\times10^{-7}$. A diagnostic that verifier rounding can zero out is reporting a
property of the stack as much as of the policy. This is why we use PCD
as the control signal.

\subsection{Testing the signal model}
\textbf{T1.} Binning per-group reward tensors by empirical success rate gives
a ZVF indicator of 1.000 at $\hat p \in \{0,1\}$ and 0.000 in every interior
bin. At run level, measured ZVF sits above $h_G(\bar p)$ at every $G$, which
is consistent with a Jensen gap: the prompt population mixes mastered and
hopeless prompts.

\textbf{T2.} Two independent grids agree. In the 368-run audit, ZVF declines
monotonically in $G$ for every model with interior success (Qwen3-32B:
0.33 $\to$ 0.18 $\to$ 0.08; Llama-3.1-8B: 0.76 $\to$ 0.65 $\to$ 0.47 across
$G = 4/8/16$). The two boundary models behave as the theory predicts:
Llama-3.2-1B, near $p\approx0$, stays at 0.98 $\to$ 0.97, and a near-saturated
Qwen3-4B model produces the single non-monotone cell. A controlled
Qwen2.5-0.5B sweep gives the same decline (Table~\ref{tab:gsweep},
Fig.~\ref{fig:gu}), with held-out accuracy essentially flat near 0.98.
Raising $G$ buys signal, but on this near-saturated task it does not buy
outcome.

\textbf{T3.} Per-group gradients are not exposed by the managed stack. On an
open backward pass with a 0.5B model and synthetic arithmetic,
$\mathrm{corr}(\lVert\nabla\rVert, \hat p(1-\hat p)) = +0.71$. We report this
as directional evidence only.

\begin{table}[t]
\caption{Group-size sweep (Qwen2.5-0.5B, 3 seeds, 40 steps): measured ZVF
against the i.i.d.\ prediction $h_G$, and gradient utilisation.}
\label{tab:gsweep}
\centering
\small
\begin{tabular}{@{}rcccc@{}}
\toprule
$G$ & ZVF (measured) & $h_G$ (theory) & Residual & GU \\
\midrule
2  & 0.838 & 0.964 & $-0.126$ & 0.162 \\
4  & 0.764 & 0.954 & $-0.191$ & 0.236 \\
8  & 0.691 & 0.923 & $-0.232$ & 0.309 \\
16 & 0.631 & 0.704 & $-0.073$ & 0.369 \\
\bottomrule
\end{tabular}
\end{table}

\begin{figure*}[t]
\centering
\includegraphics[width=0.78\textwidth]{figures/fig_gradient_utilisation.pdf}
\caption{Gradient utilisation ($1-\mathrm{ZVF}$) rises with group size with
diminishing returns (left); the relative gain per doubling of $G$ declines
(right).}
\label{fig:gu}
\end{figure*}

\subsection{Group size has a shallow optimum at best}
The marginal ZVF reduction shrinks with $G$: 0.075 from $G=2$ to 4, but only
0.059 from $G=8$ to 16. The signal-to-noise ratio reaches about 52\% of the
$\sqrt{G}$ ideal. On the near-ceiling sweep, held-out accuracy is
0.982, 0.990 and 0.978 at $G = 2, 8, 16$, which are statistically
indistinguishable at three seeds, and $G=2$ retains 100.3\% of $G=16$. A
505-task conditional utility audit prefers $G=4$ over $G=5$ by 0.00177
(95\% bootstrap CI $[0.00082, 0.00270]$): the difference is positive but small,
and it depends on the task cohort and the objective. Eq.~(\ref{eq:signal}) explains why
only a shallow optimum should exist; the measured data do not support a
universal best $G$.

\subsection{Acting on the diagnostic}
\textbf{Adaptive group size.} We audited an adaptive-$G$ controller that
escalates group size on ZVF spikes, using PCD as an aliasing guard, in a
four-arm open-trainer comparison on the same data, budget and seeds
(Table~\ref{tab:arms}). Dynamic sampling reaches perfect telemetry
(ZVF 0.00) at a 45\% rollout premium, without beating the cheaper fixed
recipe. Adaptive-$G$ ties the best held-out gain at 186 rollouts.
It is competitive with the best fixed recipe, not superior to it, and the
comparison is small-$n$ and single-task.

\begin{table}[t]
\caption{Four-arm open-trainer audit (same data, budget and seeds).}
\label{tab:arms}
\centering
\small
\begin{tabular}{@{}lccc@{}}
\toprule
Arm & Held-out gain & Mean ZVF & Rollouts \\
\midrule
GRPO & +0.500 & 0.25 & 120 \\
Dr.\ GRPO~\cite{liu2025understanding} & +0.575 & 0.27 & 120 \\
Dynamic sampling~\cite{yu2025dapo} & +0.550 & 0.00 & 174 \\
Adaptive-$G$ & +0.575 & 0.23 & 186 \\
\bottomrule
\end{tabular}
\end{table}

\textbf{Counterfactual controller trace.} Replaying 2,560 prompt-step
decisions on real reward tensors shows why escalation cannot help much here.
Of GRPO's 640 prompt-step pairs, 461 are truly degenerate (427 all-correct,
34 all-wrong) with $\hat p$ exactly 0 or 1, and at those points no $G$ restores
contrast. Bootstrapping the observed rewards to $G = 16$ restores 0 of 1,867
boundary prompts. Of 1,867 escalation events, 1,723 (92.3\%) fire on
all-correct groups and only 144 on all-wrong ones, so a symmetric ZVF trigger
spends most of its budget on prompts that are already solved. Removing those fires is
a counterfactual proposal; we did not measure whether performance is
preserved.

\textbf{Curriculum and re-baselining.} A difficulty curriculum that filters
collapsed groups drives the zero-gradient waste fraction from 0.50 to 0.00,
but at about 4.81$\times$ the sampling cost, and its held-out gain equals the
baseline's (+0.05 in both arms). At a matched 30k-token budget the two tie at
+0.028. A single-seed $G=4$ ``win'' of +0.125 did not survive multi-seed
re-testing. Re-baselining the advantages does not create learning signal;
it only turns the collapsed groups into a difficulty proxy. The
interventions that target the mechanism ZVF identifies do not improve
held-out outcomes at matched cost.

\section{Discussion and Limitations}
The theory is checked in telemetry, but its causal link to the gradient (T3)
remains toy-scale. Many managed-stack runs are single-seed early-training
snapshots of 30--50 steps, so asymptotic claims are unsupported. The sweep
task is near ceiling, where variance is compressed, so the flat retention
curve should not be generalised to hard, unsaturated tasks. The closed form
$h_G$ assumes i.i.d.\ Bernoulli rewards; measured ZVF falls below it
(Table~\ref{tab:gsweep}), which is consistent with a less extreme reward
distribution. Finally, reporting ZVF as a trajectory alongside the full
training stack is part of a broader reporting practice that we treat
separately.

\section{Conclusion}
Signal starvation is the typical state of GRPO training in our regime, and
it is driven by prompts the model already solves. A two-factor signal model
explains how group size affects it and why only a shallow optimum should
exist; two of its three predictions reproduce at scale. ZVF aliases mastery
with incapacity and is fragile under reward jitter, so PCD is the better
control signal. Neither an adaptive controller nor a curriculum converted
the diagnostic into a held-out gain at matched cost. We therefore recommend
logging ZVF and PCD as standard telemetry next to reward: a rise of more
than 0.3 sustained over ten steps is a warning. We do not recommend them as
quantities to maximise. A pre-registered bake-off of an asymmetric,
PCD-guarded controller against a static $G = 16$ is the natural next test.
